\documentclass[aspectratio=169,11pt]{beamer}
\usepackage{fontspec}
\usepackage{listings}
\usepackage{xcolor}
\usepackage{booktabs}
\usepackage{array}

% ================= Palette (identique au livre Tome 1) =================
\definecolor{PyBleu}{HTML}{1A4A7A}
\definecolor{PyBleuMoyen}{HTML}{2D7AB3}
\definecolor{PyBleuClair}{HTML}{EAF1F8}
\definecolor{PyJaune}{HTML}{FFD43B}
\definecolor{PyNuit}{HTML}{142840}
\definecolor{PyTexte}{HTML}{1F2A37}
\definecolor{PyGris}{HTML}{5A6B7E}
\definecolor{PyRouge}{HTML}{B3442D}
\definecolor{PyVert}{HTML}{2E7D5B}
\definecolor{CodeTexte}{HTML}{E8EEF6}

% ================= Code Python =================
\lstdefinestyle{py}{
  language=Python,
  keywordstyle=\color{PyJaune}\bfseries,
  stringstyle=\color{PyJaune},
  commentstyle=\color{PyGris}\itshape,
  showstringspaces=false,
  breaklines=true,
  backgroundcolor=\color{PyNuit},
  basicstyle=\ttfamily\footnotesize\color{CodeTexte},
  columns=fullflexible,
  keepspaces=true,
  xleftmargin=6pt,
  framexleftmargin=6pt,
}
\lstdefinestyle{pysmall}{
  language=Python,
  keywordstyle=\color{PyJaune}\bfseries,
  stringstyle=\color{PyJaune},
  commentstyle=\color{PyGris}\itshape,
  showstringspaces=false,
  breaklines=true,
  backgroundcolor=\color{PyNuit},
  basicstyle=\ttfamily\scriptsize\color{CodeTexte},
  columns=fullflexible,
  keepspaces=true,
  xleftmargin=5pt,
  framexleftmargin=5pt,
}

% ================= Titres de diapositives =================
\setbeamercolor{structure}{fg=PyBleu}
\setbeamerfont{frametitle}{series=\bfseries, size=\Large}
\setbeamercolor{frametitle}{bg=PyBleu, fg=white}
\setbeamertemplate{frametitle}{%
  \nointerlineskip
  \begin{beamercolorbox}[wd=\textwidth, ht=2.1ex, dp=0.9ex,
                          leftskip=0.35cm, rightskip=0.35cm]{frametitle}%
    \usebeamerfont{frametitle}\insertframetitle%
  \end{beamercolorbox}%
  \vspace*{-0.4pt}%
  {\color{PyJaune}\rule{\textwidth}{1.6pt}}%
}

% ================= Pied de page =================
\setbeamerfont{footline}{size=\tiny}
\setbeamercolor{footline}{bg=white, fg=PyGris}
\newcommand{\PiedDePage}{Python en 24 heures \textbar\ Séance 08 : Classes et encapsulation}
\setbeamertemplate{footline}{%
  \leavevmode%
  \hbox{%
    \begin{beamercolorbox}[wd=\paperwidth, ht=2.2ex, dp=1.1ex,
                            leftskip=0.35cm, rightskip=0.35cm]{footline}%
      \usebeamerfont{footline}\PiedDePage
      \hfill\insertframenumber\,/\,\inserttotalframenumber%
    \end{beamercolorbox}}%
}

\setbeamertemplate{navigation symbols}{}
\setbeamertemplate{itemize item}{\color{PyBleuMoyen}$\blacktriangleright$}
\setbeamertemplate{itemize subitem}{\color{PyBleuMoyen}\textbullet}

% ================= Blocs pedagogiques =================
\setbeamercolor{block title}{bg=PyBleu, fg=white}
\setbeamercolor{block body}{bg=PyBleuClair, fg=PyTexte}
\setbeamercolor{block title alerted}{bg=PyRouge, fg=white}
\setbeamercolor{block body alerted}{bg=PyRouge!9, fg=PyTexte}
\setbeamercolor{block title example}{bg=PyJaune, fg=PyNuit}
\setbeamercolor{block body example}{bg=PyJaune!14, fg=PyTexte}

% ================= Metadonnees =================
\hypersetup{
  pdftitle={Python en 24 heures : Séance 08 : Classes et encapsulation},
  pdfauthor={HOUNKPE Mawulikplimi Roland et DJABON Ounimborbitibou},
  pdfsubject={De l'algorithmique à la programmation orientée objet},
  pdfkeywords={Python, algorithmique, POO, formation, exercices corrigés}
}

\begin{document}

{%
\setbeamercolor{background canvas}{bg=PyNuit}
\begin{frame}[plain]
\vspace{10mm}
{\color{PyJaune}\small\bfseries FORMATION « PYTHON EN 24 HEURES » \quad\textbar\quad SÉANCE 08\par}
\vspace{4mm}
{\color{white}\Huge\bfseries Classes et objets : encapsulation\par}
\vspace{4mm}
{\color{PyBleuClair}\large Chapitres 15 et 16 du Tome 1\par}
\vspace{7mm}
{\color{PyJaune}\rule{34mm}{2.2pt}\par}
\vspace{7mm}
{\color{white}\normalsize 1 h de cours (présentation Beamer) \quad$\bullet$\quad 1 h de pratique (notebook Jupyter)\par}
\vspace{2.5mm}
{\color{PyBleuClair}\normalsize HOUNKPE Mawulikplimi Roland et DJABON Ounimborbitibou\par}
\end{frame}
}
\begin{frame}{Objectifs de la séance}
\begin{itemize}
  \item Passer du style procédural au style objet : la classe comme moule
  \item Écrire \_\_init\_\_ et comprendre le rôle de self
  \item Distinguer attributs d'instance et attributs de classe
  \item Protéger les données avec les conventions \_ et \_\_ et exposer via @property
  \item Utiliser @classmethod et @staticmethod à bon escient
\end{itemize}
\end{frame}
\begin{frame}{Rappel de la séance 07}
Jusqu'ici, les données (dict, list) et les fonctions vivaient séparément. L'objet réunit les deux dans une même entité.

\begin{itemize}
  \item LEGB : où Python cherche les noms
  \item Closures et nonlocal : une fonction qui garde son contexte
  \item Décorateurs : emballer une fonction sans la modifier
  \item Chaînes : strip, lower, split, join et le chaînage
  \item Fichiers : with open, modes r, w, a
\end{itemize}
\end{frame}
\begin{frame}[fragile]{Du procédural à l'orienté objet}
\begin{columns}[t]
\begin{column}{0.42\textwidth}
\begin{itemize}
  \setlength{\itemsep}{1.5pt}\setlength{\parskip}{0pt}
  \item Procédural : des données d'un côté, des fonctions de l'autre
  \item Objet : des données (\textbf{attributs}) et des comportements (\textbf{méthodes}) réunis
  \item La \textbf{classe} est le moule ; l'\textbf{objet} (ou instance) est le produit du moule
  \item Chaque instance garde ses propres valeurs
  \item Exemple : mille comptes bancaires, une seule classe Compte
\end{itemize}
\end{column}
\hspace{0.02\textwidth}
\begin{column}{0.55\textwidth}
\begin{lstlisting}[style=pysmall]
class Compte:
    def deposer(self, montant):
        self.solde += montant

compte_awa = Compte()
compte_awa.solde = 1000
compte_awa.deposer(500)
print(compte_awa.solde)   # 1500
\end{lstlisting}
\end{column}
\end{columns}
\end{frame}

\begin{frame}[fragile]{\_\_init\_\_ et self : la naissance d'un objet}

\begin{lstlisting}[style=py]
class Compte:
    def __init__(self, titulaire, solde=0):
        self.titulaire = titulaire
        self.solde = solde

    def deposer(self, montant):
        self.solde += montant

compte = Compte("Awa", 1000)
compte.deposer(500)
print(compte.titulaire, compte.solde)
\end{lstlisting}
\vspace{0.5pt}

{\small\color{PyGris}\itshape \_\_init\_\_ est appelé automatiquement à la création. self désigne l'instance courante : compte.deposer(500) passe compte en self.}

\end{frame}

\begin{frame}{self, mode d'emploi}
\begin{itemize}
  \item Le premier paramètre de toute méthode d'instance est \textbf{self}
  \item `compte.deposer(500)` équivaut à `Compte.deposer(compte, 500)`
  \item Oublier self dans la définition est l'erreur n°1 du débutant
  \item Les attributs naissent par affectation : `self.x = ...`
  \item On ne passe \textbf{jamais} self soi même à l'appel
\end{itemize}
\end{frame}

\begin{frame}[fragile]{Attribut d'instance contre attribut de classe}
\begin{columns}[t]
\begin{column}{0.42\textwidth}
\begin{itemize}
  \setlength{\itemsep}{1.5pt}\setlength{\parskip}{0pt}
  \item Attribut d'instance : défini dans \_\_init\_\_ via self, propre à chaque objet
  \item Attribut de classe : défini au niveau de la classe, \textbf{partagé} par toutes les instances
  \item Usage : constante commune, compteur d'instances
  \item Piège : réaffecter via self crée une copie d'instance qui \textbf{masque} l'attribut de classe
\end{itemize}
\end{column}
\hspace{0.02\textwidth}
\begin{column}{0.55\textwidth}
\begin{lstlisting}[style=pysmall]
class Compte:
    banque = "UBA"          # attribut de classe

    def __init__(self, titulaire):
        self.titulaire = titulaire  # d'instance

a = Compte("Awa")
b = Compte("Kofi")
print(a.banque, b.banque)      # UBA UBA
Compte.banque = "Ecobank"
print(a.banque, b.banque)      # Ecobank Ecobank
\end{lstlisting}
\end{column}
\end{columns}
\end{frame}

\begin{frame}{Encapsulation : protéger l'intérieur}
\begin{itemize}
  \item Principe : le monde extérieur passe par des \textbf{méthodes}, pas par les attributs bruts
  \item Convention `\_attribut` : usage interne, signal aux lecteurs
  \item Convention `\_\_attribut` : Python déforme le nom (name mangling), accès direct bloqué
  \item Objectif : valider chaque modification au lieu de subir n'importe quelle valeur
  \item Un solde négatif saisi de l'extérieur ne doit jamais être possible
\end{itemize}
\end{frame}

\begin{frame}[fragile]{@property : l'attribut intelligent}

\begin{lstlisting}[style=py]
class Compte:
    def __init__(self, titulaire, solde=0):
        self._solde = solde
    @property
    def solde(self):
        return self._solde

    @solde.setter
    def solde(self, valeur):
        if valeur < 0:
            raise ValueError("solde négatif interdit")
        self._solde = valeur
compte = Compte("Awa", 500)
compte.solde = 800      # passe par le setter
print(compte.solde)     # passe par le getter
\end{lstlisting}
\vspace{0.5pt}

{\small\color{PyGris}\itshape Le code appelant écrit compte.solde comme un attribut simple, mais chaque accès passe par vos méthodes : la validation est invisible et garantie.}

\end{frame}

\begin{frame}[fragile]{@classmethod et @staticmethod}

\begin{lstlisting}[style=py]
class Compte:
    nombre = 0
    def __init__(self):
        Compte.nombre += 1

    @classmethod
    def combien(cls):
        return cls.nombre

    @staticmethod
    def valide(montant):
        return isinstance(montant, (int, float)) and montant > 0
print(Compte.combien())   # 0
Compte(); Compte()
print(Compte.combien())   # 2
\end{lstlisting}
\vspace{0.5pt}

{\small\color{PyGris}\itshape classmethod reçoit la classe (cls) : parfait pour les fabriques et compteurs. staticmethod ne reçoit rien : une fonction rangée dans la classe par convenance.}

\end{frame}

\begin{frame}[fragile]{Fabrique via @classmethod}

\begin{lstlisting}[style=py]
class Compte:
    def __init__(self, titulaire, solde):
        self.titulaire = titulaire
        self.solde = solde

    @classmethod
    def creer_depuis_chaine(cls, texte):
        nom, montant = texte.split(":")
        return cls(nom, int(montant))

compte = Compte.creer_depuis_chaine("Awa:1500")
print(compte.titulaire, compte.solde)
\end{lstlisting}
\vspace{0.5pt}

{\small\color{PyGris}\itshape La fabrique alternative : l'appelant écrit Compte.creer\_depuis\_chaine(« Awa:1500 ») et reçoit un objet prêt. cls désigne la classe, quelles que soient ses futures filles.}

\end{frame}

\begin{frame}[fragile]{\_\_str\_\_ : un objet qui se présente}

\begin{lstlisting}[style=py]
class Compte:
    def __init__(self, titulaire, solde):
        self.titulaire = titulaire
        self.solde = solde

    def __str__(self):
        return f"Compte de {self.titulaire} : {self.solde} F"

compte = Compte("Awa", 1500)
print(compte)             # utilise __str__
\end{lstlisting}
\vspace{0.5pt}

{\small\color{PyGris}\itshape print() cherche \_\_str\_\_ ; sans lui, on voit l'adresse mémoire <\_\_main\_\_.Compte object at 0x...>. Nous approfondirons les méthodes spéciales à la séance 09.}

\end{frame}

\begin{frame}{Attention}
\begin{alertblock}{Le piège de la mutable partagée}
Un attribut de classe de type liste est partagé par toutes les instances : `Compte.historique.append(...)` écrit dans LA liste commune à tous les comptes. Règle : toute donnée propre à une instance naît dans \_\_init\_\_ via self. L'attribut de classe reste réservé aux constantes et aux compteurs.
\end{alertblock}
\end{frame}

\begin{frame}{À retenir}
\begin{block}{L'essentiel de la séance 08}
La classe définit un moule : \_\_init\_\_ fabrique l'instance et self la désigne partout. Les attributs d'instance sont individuels, les attributs de classe partagés. L'encapsulation cache les données (\_ ou \_\_) et les expose via @property avec validation. @classmethod reçoit la classe, @staticmethod n'est qu'une fonction bien rangée.
\end{block}
\end{frame}

\begin{frame}{Travaux pratiques (1 h) : Classes, objets et encapsulation}
Cinq exercices pour fabriquer vos premières classes solides. Les exercices 3 à 5 préparent le grand mini projet bancaire de la séance 10 : gardez ce code sous la main.

\begin{itemize}
  \setlength{\itemsep}{1.5pt}\setlength{\parskip}{0pt}
  \item \textbf{Exercice 1 (Basique)} : Classe Point
  \item \textbf{Exercice 2 (Basique)} : Classe Livre
  \item \textbf{Exercice 3 (Intermédiaire)} : CompteSimple validé
  \item \textbf{Exercice 4 (Intermédiaire)} : Température avec @property
  \item \textbf{Exercice 5 (Défi)} : Playlist avec compteur de créations
\end{itemize}
\vspace{4pt}

{\small\color{PyGris}Ouvrez le notebook « Seance\_08\_classes\_et\_encapsulation.ipynb » pour commencer.}
\end{frame}
\begin{frame}{Pour la prochaine séance}
\begin{exampleblock}{Conseil de travail}
Séance 09 : héritage, polymorphisme et méthodes spéciales. Vos classes vont se transmettre des comportements et apprendre à se comparer, à s'additionner et à se comporter comme des listes. Relisez \_\_str\_\_ de cette séance : il devient \_\_repr\_\_ et prend une autre dimension.
\end{exampleblock}
\end{frame}

\end{document}
